\section{讨论：把一致性推进到交互边界}
\label{sec:discussion}
\paragraph{单房间串行化是一种明确的交换。}
256 条锁带把活跃房间映射到有限同步对象，避免动态锁表的创建、回收和生命周期问题。碰撞会使无关房间暂时串行，热房间等待也会占据共享工作线程。这是当前模型为简单顺序语义付出的代价。下一步有价值的研究不是直接增加线程数，而是在热点房间与普通房间混合负载下测量排队时间；若跨房间尾延迟显著恶化，再考虑按房间就绪队列调度，让等待发生在队列中而非占用工作线程。公平性优化必须保持执行至完成入队这一原子顺序边界。

\paragraph{持久化日志不自动成为发布日志。}
当前数据库权威历史与内存完成队列各司其职。运行中请求者断开不会丢弃已提交广播，但服务进程若在提交与完成入队之间崩溃，在线通知仍可能缺失，之后通过历史补齐。若要扩展为多服务实例或要求崩溃后自动补发，需要将待发布事实与消息写入同一持久事务，并定义发布重试和消费去重。此时问题不再是扩大现有锁，而是引入持久发布日志；跨实例实时扇出也应成为显式协议能力，而不是假定 PostgreSQL 行锁能够提供它。

\paragraph{可见位置是交互事实，不是传输事实。}
显式已读将确认权留给前台用户，视口给出本次确认的上界；这一游标仍表示房间前缀，而非逐消息注意力证明。初始仅同步最近窗口、用户跳转或折叠内容时，系统不能据此推断每个历史字符都被阅读。更细粒度的可见区间或逐消息确认会改变存储成本、未读定义与 UI 负担。当前选择保留熟悉的前缀游标语义，并严格区分“收到”“同步过”和“允许显式确认”的状态。

\paragraph{边界应当成为类型与协议的一部分。}
对已转发后断开的发送返回 \code{delivery_unknown}，而对尚未入队的离线发送返回 \code{offline}，比在同一个错误字符串里解释两种情况更可靠。不可变广播缓冲区与独立写偏移、独立历史抓取游标与元数据最大值，也都是同一种设计：先使状态表达正确，再让普通控制流承担原来的特殊情况。后续演化应继续保留现有帧含义和普通 CLI 行为，在兼容字段或明确版本边界上扩展。

\section{结论}
本实验将一个原生套接字聊天系统的主要困难归结为跨边界语义的组织：会话 FIFO 不等于房间顺序，数据库提交不等于事件发布，接收最大值不等于历史进度，本地回显不等于成功发送。通过明确所有权、对齐发布门控、拆分游标以及采用真实终端旅程，系统在保持协议兼容的同时消除了多类错位。性能改进也来自同一方向：用工作就绪替代周期等待，用不可变共享字节替代重复编码，用成熟终端批处理替代逐字符网络轮询。开发文档、独立审阅和反例回归共同把这些原则转化成可以持续维护的工程结构。
